\section{Data Representation and Number Systems}
Data is internally represented as a sequence of bits (binary digits). How it is
represented depends on the data type, and certain bit strings can have
different values when taken as different data types.
\begin{definition}[Units of Data Representation]
  Each \textbf{byte} is 8 bits, each \textbf{nibble} is 4 bits, and each
  \textbf{word} is multiple bytes depending on computer architecture. For
  example, 64-bit architecture uses 8 bytes and 32-bit architecture uses 4
  bytes as the word size.
\end{definition}
$N$ bits can represent up to $2^{N}$ values. To represent $M$ values,
$\lceil\log_2 M\rceil$ bits are required. 
\subsection{Number Systems}
We consider a weighted-positional number system, in the usual \textit{base} (or
radix) 10.
\begin{definition}[Weighted-Positional Number System]
  We consider a weighted-positional number system, in \textit{base} (or \textit{radix}) R:\@
  Then we have the \textit{symbols}/\textit{digits}$=\left\{ 0,1,2,\ldots,R-1 \right\}$
  and each position has a weight of power of R.
  \[{\left( a_{n}a_{n-1}\cdots a_0.f_1f_2\cdots f_m \right)}_R=\]
  \[(a_{n}\times 10^{n})+(a_{n-1}\times 10^{n-1})+\cdots+(a_0\times 10^{0})+(f_1\times 10^{-1})+(f_2\times 10^{-2})+\cdots+(f_m\times 10^{-m})\]
\end{definition}
\begin{notation}[Number Systems Notations]
  In certain programming languages, we can specify certain bases of numbers as follows:
  \begin{multicols}{3}
    \begin{itemize}
      \item In \verb|C|:
      \begin{itemize}
        \item Prefix \verb|0| for octal
        \item Prefix \verb|0x| for hexadecimal
      \end{itemize}
      \vfill\null%
      \columnbreak%
    \item In \verb|QTSpim|:
      \begin{itemize}
        \item Prefix \verb|0x| for hexadecimal
      \end{itemize}
      \vfill\null%
      \columnbreak%
    \item In \verb|Verilog|:
      \begin{itemize}
        \item Prefix \verb|8'b| for binary
        \item Prefix \verb|8'h| for hexadecimal
        \item Prefix \verb|8'd| for decimal
      \end{itemize}
    \end{itemize}
  \end{multicols}
\end{notation}
\subsubsection{Conversion}
To convert any base-R number to decimal, we sum the symbols, with each scaled
by a weight of power of R corresponding to its position.
\begin{example}[From Base-5 to Decimal]
  \[341.24_5=3\times 5^{2}+4\times 5^{1}+1\times 5^{0}+2\times 5^{-1}+4\times 5^{-2}=96.56_{10}.\]
\end{example}
To convert a whole decimal number to any base-R number, we use
\textbf{successive division by $R$} until the quotient is 0. Each remainder
forms the answer, with the first remainder as the \textit{Least Significant
Bit} (LSB) and the last as the \textit{Most Significant Bit} (MSB).
\begin{example}[Decimal Whole Number to Binary]
  \[43_{10}=(1 \bmod 2)(2 \bmod 2)(5 \bmod 2)(10 \bmod 2)(21 \bmod 2)(43 \bmod 2)=101011_2.\] 
\end{example}
To convert a decimal fraction to any base-R number, we use \textbf{repeated
multiplication by $R$}, until the fractional product is 0 or the required
precision is achieved. The \textit{carries} produce the answer, with the first
carry as the MSB and the last as the LSB.%
\begin{example}[Decimal Fraction to Binary]
  \[0.3125_{10}=(0.3125\times 2\bmod 1)(0.625\times 2\bmod 1)(0.25 \times 2\bmod 1)(0.5\times 2\bmod 1)=0.0101.\] 
\end{example}
In general, to convert between non-decimal bases $R$ and $S$, we convert first
to decimal and then to the desired base.
\begin{remark}
  However, we may convert between binary, octal, and hexadecimal by grouping
  symbols together or expanding into multiple symbols directly.
\end{remark}

\subsection{Data Representation}
\subsubsection{Character Representation}
ASCII code and Unicode is used to represent characters. Each character is
$8$ bits in ASCII, consisting $7$ bits and $1$ parity bit. Hence characters take
ASCII code values from $0$ to $127$.
\begin{figure}[h!]
  \centering
  \includegraphics[width=0.8\textwidth]{ASCII_Table.png}
\end{figure}

\subsubsection{Integer Representation}
Numbers in C can either be unsigned or signed.\ \textit{Unsigned Numbers} can
only take non-negative values. In this case, the range of values it can take
would be from $0$ to $2^{n}-1$, where $n$ is the number of bits used.
\textit{Signed Numbers} can take both positive and negative values. However,
with the same number of bits, the range of values would hence be smaller (since
information on the sign has to be stored with a bit). There are three common
representations, namely \textbf{Sign-and-Magnitude}, \textbf{1s Complement},
and \textbf{2s Complement}. 

\paragraph{Sign-and-Magnitude}
We may represent a negative number using the MSB as a `sign bit', with $0$ for $+$ and
$1$ for $-$. We hence have, for $8$-bit integers:
In general, the range is $-2^{n-1}+1$ to $+2^{n-1}-1$ for a $n$-bit integer. We
negate a number by simply inverting the sign bit.
\[\text{Range: }11111111_2=-127_{10}\text{ to }01111111_2=+127_{10}\]
\[\text{Zeros: }00000000=+0_{10}\text{ and }10000000=-0_{10}\]

\paragraph{1s Complement}
Given a positive number $x$, its additive inverse is given in 1s-complement representation by
\[-x=2^{n}-x-1.\]
\begin{example}[$12_{10}$]
  With an $8$-bit number $00001100_2$, its negated value expressed in 1s-complement is:
  \begin{align*}
    -{\left( 00001100 \right)}_2&=(2^{8}-12-1)\\
                                &=243\\
                                &={11110011}_{1s}
  \end{align*}
\end{example}
In general, the range is $-2^{n-1}+1$ to $+2^{n-1}-1$ for a $n$-bit integer. We
negate a number by inverting all of its bits. The MSB still represents the sign
of the integer.
\[\text{Range: }10000000_{1s}=-127_{10}\text{ to }01111111_{1s}=+127_{10}\]
\[\text{Zeros: }00000000=+0_{10}\text{ and }11111111=-0_{10}\]
When performing \textbf{addition of integers in 1s complement}, add the `carry out' of
the MSB to the LSB\@. Overflow occurs if the result is opposite sign of $A$ and
$B$.

\paragraph{2s Complement}
Given a positive number $x$, its additive inverse is given in  2s-complement representation by
\[-x=2^{n}-x.\] 
\begin{example}[$12_{10}$]
  With an $8$-bit number $00001100_2$, its negated value expressed in 2s-complement is:
  \begin{align*}
    -{\left( 00001100 \right)}_2&=(2^{8}-12)\\
                                &=244\\
                                &={11110100}_{2s}
  \end{align*}
\end{example}
In general, the range is $-2^{n-1}$ to $+2^{n-1}-1$ for a $n$-bit integer. We
negate a number by inverting all of its bits, then add $1$. The MSB still represents the sign
of the integer.
\[\text{Range: }10000000_{1s}=-128_{10}\text{ to }01111111_{1s}=+127_{10}\]
\[\text{Zeros: }00000000=+0_{10}\text{ and }11111111=-0_{10}\]
When performing \textbf{addition of integers in 2s complement}, the carry out of the MSB is ignored.
Overflow occurs if the `carry in' and the `carry out' of the
MSB is different, or if the result is opposite sign of $A$ and $B$.
\begin{definition}[Sign Extension]
  To convert a $n$-bit \textbf{signed} integer to a $m$-bit signed integer
  (where $m>n$), we copy the MSB $m-n$ times to the left of the $n$-bit number.
  That is, we copy the `sign' bit of the integer.  This is
  \textit{value-preserving} only for 2s complement.
\end{definition}

\paragraph{Excess-$k$ Representation} 
Given a number in excess-$k$ representation, we may obtain its value in decimal
by taking the number as binary, and then subtracting $k$.
\begin{example}[Excess-$8$]
  For $4$-bit numbers, excess-$8$ is as shown below:

\centering
\begin{tabular}{c c c c c c c c c}
  \toprule
  Excess-$8$ Representation & $0000$ & $0001$ & $0010$ & $0011$ & $0100$ & $0101$ & $0110$ & $0111$ \\
  \midrule
                           Value & $-8$ & $-7$ & $-6$ & $-5$ & $-4$ & $-3$ & $-2$ & $-1$\\
                           \bottomrule
\end{tabular}
\begin{tabular}{c c c c c c c c c}
  \toprule
  Excess-$8$ Representation& $1000$ & $1001$ & $1010$ & $1011$ & $1100$ & $1101$ & $1110$ & $1111$\\
  \midrule
  Value & $0$ & $1$ & $2$ & $3$ & $4$ & $5$ & $6$ & $7$\\
  \bottomrule
\end{tabular}
\end{example}
\raggedright%

\subsection{Real Number Representation}
In \textbf{fixed-point representation}, we fix the number of bits allocated for
the integer and fractional parts. 
\begin{figure}[h!]
  \centering
  \includegraphics[width=0.4\textwidth]{fixedPoint.jpeg}
\end{figure}\\
We extend the limited range of fixed-point representation with
\textbf{floating-point representation}, allowing us to represent very large or
very small numbers.
\paragraph{IEEE 754 Floating-Point Representation}
The floating-point representation comprises 3 components: \textit{sign},
\textit{exponent}, and \textit{mantissa} (fraction).
\begin{itemize}
  \item Single-precision (32 bits): 1-bit sign, 8-bit exponent in excess-127, 23-bit mantissa
  \item Double-precision (64 bits): 1-bit sign, 11-bit exponent in excess-1023, 52-bit mantissa
\end{itemize}
The sign bit has $0$ for $+$, and $1$ for $-$. The exponent has an implicit
base $2$, and is stored in excess representation, whereas the mantissa is normalised
with an implicit leading bit 1. That is, the mantissa only stores the bits
after the binary point.
\[v={\left( -1 \right)} ^{sign}\times 1.fraction\times 2^{exponent-bias}.\] 
\begin{example}[$-6.5_{10}$]
  To represent a given value, we first convert into binary, and normalise the mantissa:
  \[-6.5_{10}=-110.1_2=-1.101_2\times 2^{2}.\]
  We hence have the sign bit $=1$, exponent
  $=2_{10}=10000001_{\text{excess-}127}$, mantissa $=10100000000000000000000$.
\end{example}
\begin{definition}[Underflow and Overflow]
  With the discrete representation of numbers, we have numbers that are
  non-representable on the number line, as shown:
  \begin{figure}[h!]
    \centering
    \includegraphics[width=0.5\textwidth]{floatingUnderOverflow.png}
  \end{figure}
  As such, associativity is not preserved, since 
  \[\left( A \cdot B \right) \cdot C=round{\left( round{\left( A \cdot B \right)} \cdot C \right) }.\]
\end{definition}
Depending on implementation, rounding can be to the nearest (default), towards
zero, towards positive infinity, or negative infinity. The operation is done
with 3 extra bits (guard, round, sticky), and then rounded to the nearest
representable number.
